\chapter{One Circuit, Four Views}
\label{ch:fourviews}
Part~II introduced differential equations, state-space models, transfer functions,
frequency responses, and feedback. Chapters~7--9 supplied circuit components,
conservation laws, and phasor impedance. This chapter is the bridge: it connects a
physical circuit to the generic system model that describes its dynamics, so each
later circuit chapter can use both languages without starting over.

\section{The Four Views}
A circuit can be understood through:
\begin{enumerate}
\item a \textbf{physical network} governed by conservation laws (KCL, KVL) and
element laws;
\item a \textbf{differential equation} in the natural energy-storage variables;
\item a \textbf{state-space model}, whose eigenvalues equal the transfer-function
poles when the selected realization is minimal;
\item a \textbf{phasor / frequency-response model}, the transfer function on the
\(\jj\omega\) axis.
\end{enumerate}
\cite[pp.~199--203, 249--266]{leigh2012control}
A fifth lens is \textbf{energy} storage and dissipation. The views are
complementary: each makes a different feature of
the circuit obvious, and together they describe one consistent object. When a
quick cross-check between two of them disagrees, a sign convention, a reference
direction, or an assumption needs attention.

\section{The Recipe for Part III}
The next chapters each apply the following recipe to specific circuits:
\begin{enumerate}
\item draw the schematic and name the state variables (capacitor voltages,
inductor currents);
\item apply KVL/KCL and the element laws to get the differential equation;
\item write the state-space model, distinguish its eigenvalues from poles visible
at the selected input/output, and identify when the realization is minimal;
\item match the characteristic denominator to the generic first- or second-order
system of \cref{ch:linsys}, identifying \(\tau\), or \(\omega_0\) and \(\zeta\);
then derive the numerator for the chosen source and output;
\item interpret the frequency response (Bode) and the pole motion (root locus);
\item close with the energy balance and a worked example.
\end{enumerate}

\section{The Map from Circuits to Systems}
Each independent capacitor voltage or inductor current contributes one state;
resistors dissipate energy but do not raise the order. Thus a circuit with one
independent energy store is first order, and one with two is second order. Four
circuits carry most of the exam's analog content:

The table maps each circuit's \emph{characteristic dynamics}. A complete transfer
function also requires a specified source and output; those choices set the
numerator and forcing term.

\begin{center}
\begin{tabularx}{\textwidth}{@{}L{0.15\textwidth}L{0.19\textwidth}L{0.24\textwidth}Y@{}}
\toprule
Circuit & Order / state(s) & Generic system & Key parameters \\
\midrule
Series RC & first / \(v_C\) & first-order characteristic denominator & \(\tau=RC\), pole \(-1/RC\) \\
Series RL & first / \(i_L\) & first-order characteristic denominator & \(\tau=L/R\), pole \(-R/L\) \\
Series RLC & second / \(v_C,i_L\) & second-order characteristic denominator & \(\omega_0=\tfrac{1}{\sqrt{LC}}\), \(\zeta=\tfrac{R}{2}\sqrt{C/L}\) \\
Parallel RLC & second / \(v,i_L\) & second-order characteristic denominator & \(\omega_0=\tfrac{1}{\sqrt{LC}}\), \(\zeta=\tfrac{1}{2R}\sqrt{L/C}\) \\
\bottomrule
\end{tabularx}
\end{center}

Once a circuit is matched to its characteristic row, the Part~II time,
pole-geometry, and frequency results can be specialized to the selected input and
output. After topology and algebraic constraints are accounted for, additional
independent capacitor voltages and inductor currents increase physical state order.
A particular transfer function can expose fewer poles if modes are uncontrollable
or unobservable at that port
\cite[pp.~199--203, 264--266]{leigh2012control}.

For stable, minimal, all-pole higher-order lumped circuits, the factorization in
\cref{ch:highorder} organizes those states into first- and second-order sections.
A transmission line is instead a distributed model, developed in \cref{ch:lines}.
